\scanpage{113}
\noindent\(\left\{\vcenter{\hbox{\begin{minipage}{.91\linewidth}
\uline{Théorème} Soit \deleted{\unsure{un}} $k$ anneau, $\alpha\in k$, d'où $u_\alpha:k^{2}\to k^2$
\deleted{transformation} \inserted{\unsure{autom.|automorph.}} affine
\[
u_\alpha(x,y)=(1-y+\alpha x,x)
\]
Satisfaisant, \deleted{\unsure{si}} posant $s_n=u_\alpha^n(s_0)$ (où $s_0=(0,0)$)
\[
\begin{aligned}
\deleted{\unsure{u_\alpha^n(s_0)|u_\alpha(s_0)}}\qquad s_1&=u_\alpha^{\deleted{\unsure{n}}}(s_0)=(1,0)\\
s_{-1}&=u_\alpha^{-1}(s_0)=(0,1)\quad\text{i.e.}\quad u_\alpha(0,1)=(0,0).
\end{aligned}
\]
Soit $n\in\mathbb Z$. Conditions équivalentes

a) $u_\alpha^n=\mathrm{id}$

b) $s_n=s_0$ i.e $u_\alpha^ns_0=s_0$

b$_{\unsure{i|i+}}$) (si $i\in\mathbb Z$ fixé) $u_\alpha^ns_i=s_i$

d) $F_n(\alpha)=0$
\end{minipage}}}\right.\)

a $\Rightarrow$ b trivial

b $\Rightarrow$ a car b) implique \inserted{(en multipliant par $u_\alpha^i$)} $u_\alpha^ns_i=s_i$ $\forall i$, en
partic si $i=\pm1$ on trouve $\begin{cases}u_\alpha^n(s_{-1})=s_{-1}\\u_\alpha^n(s_1)=s_1\end{cases}$
d'où $u_\alpha^n=\mathrm{id}$
\[
\begin{aligned}
\text{b)}&\Rightarrow\text{b}_i\text{)}\quad\text{en multipliant par }u_\alpha^i\\
\text{b}_i\text{)}&\Rightarrow\text{b)}\qquad\rule{2.8cm}{.3pt}\qquad u_\alpha^{-i}
\end{aligned}
\]
donc $a\Longleftrightarrow b\Longleftrightarrow c$

Comme $s_n=(F_n(\alpha)F_{n+1}(\alpha),F_{n-1}(\alpha)F_n(\alpha))$, on voit
que d) $\Rightarrow$ b) i.e. $F_n(\alpha)=0\Rightarrow s_n=s_0\ (=(0,0))$
Inversement, si $F_n(\alpha)F_{n+1}(\alpha)=0$, $F_n(\alpha)F_{n-1}(\alpha)=0$ 
il s'ensuit $F_n(\alpha)(a F_{n-1}(\alpha)+b F_{n+1}(\alpha))=0$ $\forall a,b\in k$.
Or $\exists A_n(U),B_n(U)\in\mathbb Z[U]$ tels que
\[
A_n(U)F_{n-1}(U)+B_n(U)F_{n+1}(U)=1,
\]
d'où, faisant $a=A_n(\alpha)$, $b=B_n(\alpha)$, la relation $F_n(\alpha)=1$,
cqfd\ldots
% archive sheet number: 112
